\subsection{Application Profiling and Hardware/Software Partitioning}
The second stage determines which application functions should remain on the CPU and which should be implemented as custom accelerators. Hotspot detection alone cannot make this decision: a frequently executed function may provide little net benefit after accelerator latency, hardware cost, invocation overhead, and data movement are considered. SoCPilot therefore treats profiling as candidate extraction and defers the mapping decision to an explicit hardware/software partitioning model.

\textbf{Application profiling and candidate extraction.}
SoCPilot compiles and executes the application with representative inputs and uses Callgrind to collect function-level instruction counts, call frequencies, and calling relationships. System routines, library wrappers, initialization code, and error-handling paths are removed from the profile, leaving application-defined functions that account for substantial execution time. A lightweight static pass then extracts each candidate with its type definitions and calling context and records potential synthesis obstacles, including recursion, dynamic allocation, irregular pointer accesses, and complex control flow. This pass determines whether a function can be evaluated by the HLS flow; it does not decide whether the function should be mapped to hardware. The LLM combines the dynamic profile and static call information into a task graph $G=(V,E)$, where nodes represent profiled functions and edges record control and data dependencies.

\textbf{LLM-assisted hardware evaluation.}
For each candidate hotspot, SoCPilot reuses the LLM-guided HLS generation and optimization procedure introduced in HLSPilot \cite{xiong2024hlspilot}. The LLM isolates or refactors the candidate into synthesizable C/C++, generates a functional testbench, and iteratively inserts HLS pragmas for pipelining, loop unrolling, and array partitioning. HLS reports provide the latency of one accelerator invocation and the synthesis information used by the common analytical area model to obtain an accelerator-area estimate $a_i$. The per-invocation latency is multiplied by the profiled call count to obtain the workload-level accelerator time $t_i^{\mathrm{acc}}$. Candidates that fail synthesis or verification remain software tasks; candidates with valid implementations proceed to the partitioning algorithm, even when their estimated area is high. Thus, HLS compatibility, latency, and area cost are measured or tool-derived rather than inferred from source appearance alone.

\textbf{CPU and communication models.}
The workload-level CPU latency $t_i^{\mathrm{cpu}}$ of each task is obtained from profiling and fast simulation with Sniper \cite{sniper}. For a hardware task, the model also accumulates accelerator launch and synchronization overhead as $t_i^{\mathrm{launch}}$. Each dependency $(i,j)\in E$ is annotated with its aggregate transfer volume $D_{ij}$ over the representative execution. If two dependent tasks are placed in different execution domains, their communication cost is approximated as
\begin{equation}
c_{ij}(x_i,x_j)=
\mathbb{I}[x_i\neq x_j]\left(L_{\mathrm{if}}+\frac{D_{ij}}{B_{\mathrm{if}}}\right),
\end{equation}
where $x_i\in\{0,1\}$ denotes software or hardware execution, and $L_{\mathrm{if}}$ and $B_{\mathrm{if}}$ are the interface latency and effective bandwidth. The execution cost of task $i$ is
\begin{equation}
d_i(x_i)=(1-x_i)t_i^{\mathrm{cpu}}+
x_i\left(t_i^{\mathrm{acc}}+t_i^{\mathrm{launch}}\right).
\end{equation}

\textbf{Partitioning formulation and algorithm.}
Given these estimates, SoCPilot minimizes the predicted completion time of the task graph subject to the application-specific area budget:
\begin{equation}
\begin{aligned}
\min_{\mathbf{x}}\quad &T_{\mathrm{app}}(\mathbf{x})
=\operatorname{Makespan}\!\left(G,\{d_i\},\{c_{ij}\}\right),\\
\mathrm{s.t.}\quad &\sum_{i\in V}x_i a_i\leq A^{\max}.
\end{aligned}
\label{eq:partition}
\end{equation}
Non-candidate functions are fixed to software. The implementation uses the marginal-gain heuristic in Algorithm~\ref{alg:hw-sw-partitioning}. Starting from an all-software mapping, it tentatively maps each feasible candidate to hardware, recomputes the task-graph completion time including boundary communication, and selects the candidate with the largest latency reduction per normalized unit of estimated area. The process stops when no feasible candidate produces a positive end-to-end benefit.

\begin{algorithm}
\caption{Hardware/Software Partitioning}
\label{alg:hw-sw-partitioning}
\KwIn{Task graph $G=(V,E)$; candidates $C$; timing estimates; area estimates $a_i$; area budget $A^{\max}$}
\KwOut{Hardware set $T_H$ and software set $T_S$}

$x_i\leftarrow 0,\ \forall i\in V$; $A_{\mathrm{used}}\leftarrow 0$\;
\While{true}{
    $C_f\leftarrow\{i\in C\mid x_i=0,\ A_{\mathrm{used}}+a_i\leq A^{\max}\}$\;
    \If{$C_f=\emptyset$}{break\;}
    \ForEach{$i\in C_f$}{
        $\mathbf{x}^{(i)}\leftarrow\mathbf{x}$; $x_i^{(i)}\leftarrow 1$\;
        $\Delta_i\leftarrow T_{\mathrm{app}}(\mathbf{x})
        -T_{\mathrm{app}}(\mathbf{x}^{(i)})$\;
        $\rho_i\leftarrow a_i/A^{\max}$\;
        $s_i\leftarrow\Delta_i/\rho_i$\;
    }
    $i^\star\leftarrow\arg\max_{i\in C_f}s_i$\;
    \If{$\Delta_{i^\star}\leq 0$}{break\;}
    $x_{i^\star}\leftarrow 1$\;
    $A_{\mathrm{used}}\leftarrow A_{\mathrm{used}}+a_{i^\star}$\;
}
$T_H\leftarrow\{i\mid x_i=1\}$\;
$T_S\leftarrow V\setminus T_H$\;
\Return{$T_H,T_S$}
\end{algorithm}

The resulting partition fixes the accelerator set and provides the initial architecture for SoC generation. The subsequent DSE stage fine-tunes processor, memory-system, interconnect, and accelerator parameters without repeatedly changing this mapping. Only when no feasible design can be found under the specified constraints does SoCPilot return to the partitioning stage and revise the hardware task set.

